\section{The process on the infinite lattice}\label{sec:infinite}

We now give the construction and prove
Theorem~\ref{thm:infinite}. No global chain of all events is ordered
on the infinite lattice; there are infinitely many marks in any
positive time interval. Local chronological dependency estimates
provide the appropriate construction.

\subsection{Graphical existence}

\begin{lemma}[Finite influence on bounded intervals]\label{lem:graphical}
For independent rate-one correction clocks and independent sector
noise clocks of rate at most \(2\eta\) on each plaquette of \(\Z^4\),
finite-volume constructions started empty agree on any fixed finite
target set, throughout each bounded time interval, for all sufficiently
large cuts, almost surely. Their common limit is a local
right-continuous process with left limits. It preserves even degree
at every syndrome vertex.
\end{lemma}
\begin{proof}
Superpose the correction and noise marks of a face. Its rate is at
most \(\Lambda=1+2\eta\), uniformly over faces. Fix a finite target
set \(F\) and time length \(H_*>0\). If a discrepancy at a remote
boundary affects an edge of \(F\), tracing dependencies backwards
through the local maps produces a chronological chain of intersecting
marked faces ending next to \(F\). A correction reads only its two
tested edges and toggles its boundary, so any discrepancy it creates
has such a predecessor. A noise mark acts by the same fixed toggle
in both constructions and cannot create a discrepancy from equal
inputs. Allowing it as a possible dependency step only enlarges
the set of chains being counted.

There are at most \(C_F\) terminal faces next to \(F\), and at most
96 predecessor face labels for each subsequent step. For any
specified sequence of \(m\) labels, the ordered factorial-moment
integral is bounded by \((\Lambda H_*)^m/m!\), including repeated
labels. Therefore
\begin{equation}\label{eq:infinite-influence}
 \Prb\{\text{such a chain of length }m\}
 \le C_F\frac{(96\Lambda H_*)^m}{m!}.
\end{equation}
The series of the right sides over \(m\) converges. A dependency
from distance tending to infinity requires length tending to
infinity, because consecutive face anchors differ by at most one
in each coordinate. By the summability in
\eqref{eq:infinite-influence}, almost surely sufficiently long
chains do not occur. Hence all sufficiently large finite cuts
agree on \(F\) throughout \([0,H_*]\). Any boundary convention
confined to the cut boundary has the same limit.

Intersect these probability-one events over increasing finite target
boxes and integer \(H_*\). The resulting limits agree on overlaps
and define consistent local paths for all \(t\ge0\).
Each fixed edge is incident to finitely many face processes and
is therefore touched only finitely often on a bounded interval.
The paths are right-continuous with left limits. Every vertex
initially has degree zero, and each local face toggle changes its
degree by zero modulo two. This parity property passes from the
finite constructions to the local limit.
\end{proof}

The same construction can be made simultaneously in the two
marginal sectors with shared \(Y\) marks. The assertions here need
only the construction and law of one named sector. They do not
require a definition of an infinite-volume encoded quantum state.

\subsection{Uniform component tails}

Use ordinary blocks
\[
 Q_z=Bz+\{0,\ldots,B-1\}^4,\qquad z\in\Z^4,
\]
with the fixed \(B,T,R\) already chosen for a sufficiently small
noise threshold. The local tests and clean cuts of
Section~\ref{sec:local} are finite objects, so they apply to the
graphical limit. Their probabilities and predecessor implications
are unchanged. By induction, physical visits satisfy
\eqref{eq:domination} for every finite slab index.

Every comparison root at finite date has a finite dependency cone.
There is no periodic identification and no bound on its permissible
depth. Thus Proposition~\ref{prop:witness} gives the certificates
uniformly in the root date. Proposition~\ref{prop:path} applies on
the entire half-space of dates \(k>0\). One can also obtain its
half-space assertion by applying the finite-path argument directly:
each finite witnessing path has a finite maximum date, and the
deterministic chain count already includes all possible center
dates with the stated distance restrictions.

\begin{proof}[Proof of Theorem~\ref{thm:infinite}]
Existence and parity follow from Lemma~\ref{lem:graphical}.
Fix \(x,t,d\), and let \(z_x\) be the block containing \(x\).
Take the slab containing \(t\), with index \(k\), as specified after
the theorem. Suppose at a common instant of that slab \(x\) belongs
to a component of diameter at least \(d\).
For a finite component, the triangle inequality implies that
some vertex has distance at least \(d/2\) from \(x\); if the
diameter is only a supremum one can decrease this target radius
arbitrarily and pass to a limit of the corresponding events.
An unbounded component likewise supplies a finite connecting
edge path to any fixed radius.

The blocks along a connecting physical edge path form a spatial
path in \(\Proxy^0(\cdot,k)\) by \eqref{eq:domination}. Starting at
\(z_x\), it reaches coarse distance at least
\(d/(2B)-1\). Proposition~\ref{prop:path} bounds the probability
by
\[
 C\exp\left[-\kappa\max\{0,d/(2B)-1\}\right]
 \le C'e^{-\mu d}
\]
for \(\mu=\kappa/(2B)\), after enlarging \(C'\).
This proves the whole-slab estimate, uniformly in \(x\) and \(k\).
The event at the specified time \(t\) is a subset. At \(t=0\) the
instantaneous syndrome is empty; the whole-slab estimate uses
the slab beginning there. These statements prove
\eqref{eq:infinite-main} with uniform constants.

For the low-noise density limit it is useful that the cost \(A\)
in Proposition~\ref{prop:witness} is adjustable. At any specified
block and date,
\[
 \Prb(\Proxy^0(z,k))\le\sum_{r\ge1}e^{-Ar}
                         =\frac{1}{e^A-1}.
\]
Given \(\epsilon>0\), choose \(A\) large enough to make this
quantity less than \(\epsilon\), and large enough for the path
estimate if desired. Next choose \(q\), then \(B\), and then a
positive threshold \(\eta_\epsilon\) by the certificate proof
and Lemma~\ref{lem:bad-prob}. For every family of rates bounded
by \(\eta_\epsilon\), occupation of \(x\) at any deterministic
time implies this block proxy event in its slab. Hence its
probability is at most \(\epsilon\), uniformly in \(x,t\) and
the rate family. This is precisely the stated limit. The block
scale in this density argument may depend on \(\epsilon\);
the argument asserts no uniform limiting value for the decay
constant \(\mu\) as the threshold is varied.
\end{proof}

For each specified vertex and time the diameter tail also rules out
membership in an infinite component by letting \(d\) increase.
This is a rooted estimate and carries no factor for the volume of
the lattice. At a fixed positive noise rate new local syndromes
continue to be created, so uniform-in-time tails do not mean that
the entire infinite syndrome eventually becomes empty.
